% GENERATED FILE. Edit canonical lessons, then run npm run generate:books.
% canonical-source-hash: fnv1a64:e25c707b0e985177
% canonical-lessons: PT-C294-investigacao, PT-C294-investigar, PT-C294-inspetor, PT-C294-detetive, PT-C294-indicio

\chapter{A investigação}
\label{ch:pt-a-investigacao}
\hlchaptermodality{294}

\begin{chapteropening}
\textbf{By the end of this chapter:} \emph{I can talk about a police investigation: the inspector, the detective and the clues.}

Complete the last lesson of chapter 294: I can talk about a police investigation: the inspector, the detective and the clues.
\end{chapteropening}

\section[a investigação]{a investigação — the investigation; research}

\label{lesson:PT-C294-investigacao}



\begin{quote}
Before the new one: say the Portuguese for to keep watch over, then the Portuguese for surveillance.
\end{quote}

\subsection*{You'll want to know: a investigação}
\textbf{a investigação} — \textquotedblleft{}the investigation; research\textquotedblright{}.

\begin{quote}
A investigação durou dois anos.
\end{quote}

\begin{quote}
\emph{The investigation lasted two years.}
\end{quote}

At a university it means research.

\begin{cousinweb}
\emph{investigação} is built on \emph{investigar}, from Latin \emph{investīgāre} \textquotedblleft{}to track down\textquotedblright{} (EtymDB), from \emph{vestīgium} \textquotedblleft{}a footprint\textquotedblright{}.
\end{cousinweb}

\begin{grammarlens}[title={What you've built}]
One more word for telling, explaining and describing at B1.
\end{grammarlens}

\subsection*{Your turn}
\begin{itemize}
\raggedright
  \item \emph{Say it:} \emph{a investigação}
  \item \emph{Say it:} \emph{a investigação}, once more
  \item \emph{Say it:} a sentence of your own with \emph{a investigação}
  \item \emph{Recall:} say the Portuguese for to keep watch over, then the Portuguese for surveillance, then say \emph{a investigação} again
\end{itemize}

\begin{tcolorbox}[breakable,colback=teal!4,colframe=teal!35!black,title={Before you move on}]
What does a investigação mean? (\textquotedblleft{}The investigation; research\textquotedblright{}.) Say it once more.
\end{tcolorbox}
\section[investigar]{investigar — to investigate}

\label{lesson:PT-C294-investigar}



\begin{quote}
Before the new one: say the Portuguese for surveillance, then the Portuguese for the investigation.
\end{quote}

\subsection*{You'll want to know: investigar}
\textbf{investigar} — \textquotedblleft{}to investigate\textquotedblright{}.

\begin{quote}
A polícia está a investigar o homicídio.
\end{quote}

\begin{quote}
\emph{The police are investigating the killing.}
\end{quote}

\emph{Estar a investigar} in Portugal; in Brazil, \emph{estar investigando}.

\begin{cousinweb}
\emph{investigar} is a learned borrowing from Latin \emph{investīgāre} \textquotedblleft{}to track down, to search out\textquotedblright{} (EtymDB), built on \emph{vestīgium} \textquotedblleft{}a footprint\textquotedblright{}.
\end{cousinweb}

\begin{grammarlens}[title={What you've built}]
One more word for telling, explaining and describing at B1.
\end{grammarlens}

\subsection*{Your turn}
\begin{itemize}
\raggedright
  \item \emph{Say it:} \emph{investigar}
  \item \emph{Say it:} \emph{investigar}, once more
  \item \emph{Say it:} a sentence of your own with \emph{investigar}
  \item \emph{Recall:} say the Portuguese for surveillance, then the Portuguese for the investigation, then say \emph{investigar} again
\end{itemize}

\begin{tcolorbox}[breakable,colback=teal!4,colframe=teal!35!black,title={Before you move on}]
What does investigar mean? (\textquotedblleft{}To investigate\textquotedblright{}.) Say it once more.
\end{tcolorbox}
\section[o inspetor]{o inspetor — the inspector}

\label{lesson:PT-C294-inspetor}



\begin{quote}
Before the new one: say the Portuguese for the investigation, then the Portuguese for to investigate.
\end{quote}

\subsection*{You'll want to know: o inspetor}
\textbf{o inspetor} — \textquotedblleft{}the inspector\textquotedblright{}.

\begin{quote}
O inspetor fez muitas perguntas.
\end{quote}

\begin{quote}
\emph{The inspector asked a lot of questions.}
\end{quote}

A woman is \emph{a inspetora}. Before the 1990 spelling reform, \emph{inspector}.

\begin{cousinweb}
\emph{inspetor} is a learned borrowing from Latin \emph{īnspector} \textquotedblleft{}one who looks into\textquotedblright{}, from \emph{īnspicere} \textquotedblleft{}to look into\textquotedblright{} (Lewis \& Short).
\end{cousinweb}

\begin{grammarlens}[title={What you've built}]
One more word for telling, explaining and describing at B1.
\end{grammarlens}

\subsection*{Your turn}
\begin{itemize}
\raggedright
  \item \emph{Say it:} \emph{o inspetor}
  \item \emph{Say it:} \emph{o inspetor}, once more
  \item \emph{Say it:} a sentence of your own with \emph{o inspetor}
  \item \emph{Recall:} say the Portuguese for the investigation, then the Portuguese for to investigate, then say \emph{o inspetor} again
\end{itemize}

\begin{tcolorbox}[breakable,colback=teal!4,colframe=teal!35!black,title={Before you move on}]
What does o inspetor mean? (\textquotedblleft{}The inspector\textquotedblright{}.) Say it once more.
\end{tcolorbox}
\section[o detetive]{o detetive — the detective}

\label{lesson:PT-C294-detetive}



\begin{quote}
Before the new one: say the Portuguese for to investigate, then the Portuguese for the inspector.
\end{quote}

\subsection*{You'll want to know: o detetive}
\textbf{o detetive} — \textquotedblleft{}the detective\textquotedblright{}.

\begin{quote}
O detetive seguiu o homem até casa.
\end{quote}

\begin{quote}
\emph{The detective followed the man home.}
\end{quote}

\emph{Um romance policial} is a detective novel.

\begin{cousinweb}
\emph{detetive} was borrowed from English \emph{detective} (EtymDB), a word that goes back to Latin \emph{dētegere} \textquotedblleft{}to uncover\textquotedblright{} (Lewis \& Short).
\end{cousinweb}

\begin{grammarlens}[title={What you've built}]
One more word for telling, explaining and describing at B1.
\end{grammarlens}

\subsection*{Your turn}
\begin{itemize}
\raggedright
  \item \emph{Say it:} \emph{o detetive}
  \item \emph{Say it:} \emph{o detetive}, once more
  \item \emph{Say it:} a sentence of your own with \emph{o detetive}
  \item \emph{Recall:} say the Portuguese for to investigate, then the Portuguese for the inspector, then say \emph{o detetive} again
\end{itemize}

\begin{tcolorbox}[breakable,colback=teal!4,colframe=teal!35!black,title={Before you move on}]
What does o detetive mean? (\textquotedblleft{}The detective\textquotedblright{}.) Say it once more.
\end{tcolorbox}
\section[o indício]{o indício — the clue, the sign}

\label{lesson:PT-C294-indicio}



\begin{quote}
Before the new one: say the Portuguese for the inspector, then the Portuguese for the detective.
\end{quote}

\subsection*{You'll want to know: o indício}
\textbf{o indício} — \textquotedblleft{}the clue, the sign\textquotedblright{}.

\begin{quote}
Não encontraram nenhum indício de assalto.
\end{quote}

\begin{quote}
\emph{They found no sign of a robbery.}
\end{quote}

Something that points to what happened.

\begin{cousinweb}
\emph{indício} is a learned borrowing from Latin \emph{indicium} \textquotedblleft{}a sign, evidence\textquotedblright{}, from \emph{index} \textquotedblleft{}one who points out\textquotedblright{} (Lewis \& Short).
\end{cousinweb}

\begin{grammarlens}[title={What you've built}]
One more word for telling, explaining and describing at B1.
\end{grammarlens}

\subsection*{Your turn}
\begin{itemize}
\raggedright
  \item \emph{Say it:} \emph{o indício}
  \item \emph{Say it:} \emph{o indício}, once more
  \item \emph{Say it:} a sentence of your own with \emph{o indício}
  \item \emph{Recall:} say the Portuguese for the inspector, then the Portuguese for the detective, then say \emph{o indício} again
\end{itemize}

\begin{tcolorbox}[breakable,colback=teal!4,colframe=teal!35!black,title={Before you move on}]
What does o indício mean? (\textquotedblleft{}The clue, the sign\textquotedblright{}.) Say it once more.
\end{tcolorbox}
